\section{La mise en \oe{}uvre sociale~: Diversité, Égalité et Inclusion}

\subsection{Gouvernance et ambition du DEI}

La politique \glos{DEI} (Diversity, Equity and Inclusion) s'appuie sur un dispositif de gouvernance solide. Le \textbf{comité DEI} est présidé par le PDG, avec la \textbf{direction exécutive (SLT)} en rôle de conseil et la \textbf{direction des ressources humaines} en charge de la vision opérationnelle. Des \textit{ambassadeurs DEI} contribuent aux bonnes pratiques \cite{iveco_annual_2025}. En 2025, six priorités ont été réaffirmées~: alignement stratégique et gouvernance, culture et sensibilisation, attraction, développement et rétention, formation et capacités, engagement et communication, mesure et comparaisons externes \cite{iveco_annual_2025}.

\subsection{Égalité de genre et rémunération}

Les femmes représentent \textbf{20,5\,\%} de l'effectif mondial d'Iveco Group. Au niveau de la direction exécutive (SLT), cette part atteint \textbf{31\,\%} en 2025, après avoir atteint le seuil d'un tiers en 2024 \cite{iveco_annual_2025}. Un objectif stratégique a été inscrit au plan d'affaires : atteindre \textbf{30\,\% de femmes dans les fonctions de bureau} d'ici fin 2028 \cite{iveco_dei_2025}.

Le groupe a obtenu la certification \glos{EDGE} Assess en 2024, deux ans avant l'échéance prévue, couvrant \textbf{90\,\%} de l'effectif \cite{iveco_annual_2025}. L'évaluation a porté sur l'équité de rémunération, la représentation féminine et les politiques organisationnelles. Par ailleurs, l'écart de rémunération entre les sexes est situé dans les seuils de la directive européenne sur l'équité de rémunération \cite{iveco_annual_2025}.

\subsection{Inclusion du handicap et diversité culturelle}

En 2025, une évaluation portant sur l'emploi de personnes en situation de handicap a été conduite dans \textbf{21 pays} couvrant plus de \textbf{93,5\,\%} des effectifs. Les personnes en situation de handicap représentent \textbf{3,6\,\%} de l'effectif mondial recensé, dont \textbf{21,2\,\%} de femmes et \textbf{78,8\,\%} d'hommes \cite{iveco_annual_2025}.

Une étude sur la nationalité a également été menée dans neuf pays (83\,\% de l'effectif) : \textbf{7,7\,\%} des salariés sont d'une nationalité différente de celle du pays d'accueil, avec une part plus élevée chez les femmes (8,9\,\%) que chez les hommes (7,4\,\%). Les principaux pays d'accueil sont l'Allemagne (18\,\%), les États-Unis (11\,\%) et l'Italie (10\,\%) \cite{iveco_annual_2025}.

Ces chiffres confirment une volonté de mesurer objectivement la diversité. Toutefois, les données relatives au handicap restent partielles dans les pays ne prévoyant pas d'obligation légale, car elles reposent sur la libre déclaration des salariés.

\subsection{Santé, sécurité et bien-être au travail}

La politique Santé et Sécurité d'Iveco Group s'applique à l'ensemble des travailleurs, y compris les sous-traitants et intérimaires. Le système de management est certifié \glos{ISO} 45001 sur \textbf{27 sites}, couvrant 25\,206 salariés (72\,\% des effectifs du périmètre) \cite{iveco_annual_2025}. En 2025, le groupe a réalisé \textbf{298 audits internes} (28\,876 salariés, soit 82\,\%) et \textbf{27 audits externes} (25\,206 salariés, soit 72\,\%) \cite{iveco_annual_2025}.

Le Tableau~\ref{tab:hse} synthétise les principaux indicateurs.

\begin{table}[ht]
\centering
\caption[Indicateurs Santé et Sécurité 2025]{Indicateurs Santé et Sécurité 2025 (source~: \cite{iveco_annual_2025})}
\label{tab:hse}
\footnotesize
\begin{tabularx}{\textwidth}{@{}lr>{\raggedright\arraybackslash}X@{}}
\toprule
\textbf{Indicateur} & \textbf{Valeur} & \textbf{Commentaire} \\
\midrule
Heures de formation SST & 182\,832\,h & dont 126\,011\,h en poste \\
Salariés formés en SST & 23\,002 & 65\,\% de l'effectif (81\,\% ouvriers) \\
Taux de fréquence de blessures & 1,445 & -17,4\,\% par rapport à 2024~; -51\,\% par rapport à 2019 \\
Blessures avec arrêt & 86 & dont 100\,\% liées au travail \\
Heures travaillées & 59\,500\,256\,h & 99,7\,\% des effectifs \\
Quasi-accidents signalés & 1\,498 & actions correctives déployées \\
Budget SST & 56,6\,M\texteuro & 51,5\,M\texteuro{} sécurité et 5,1\,M\texteuro{} santé \\
\bottomrule
\end{tabularx}
\end{table}

\begin{figure}[htbp]
\centering
\begin{tikzpicture}
\begin{axis}[
    width=0.85\textwidth,
    height=0.35\textwidth,
    xlabel={Année},
    ylabel={Taux de blessures par million d'heures},
    xmin=2018.5, xmax=2026.5,
    ymin=0, ymax=3.2,
    xtick={2019,2024,2025,2026},
    xticklabels={2019,2024,2025,2026 (cible)},
    legend pos=north east,
    grid=major,
    legend style={font=\small},
]
\addplot[color=ivecoblue, mark=*, thick] coordinates {(2019,2.951) (2024,1.75) (2025,1.445)};
\addlegendentry{Taux constaté}
\addplot[color=ivecoorange, mark=triangle*, dashed, thick] coordinates {(2026,1.771)};
\addlegendentry{Cible 2026 (-40\,\%)}
\end{axis}
\end{tikzpicture}
\caption{Évolution du taux de fréquence de blessures par million d'heures (2019-2025) et cible 2026. La baisse de 51\,\% entre 2019 et 2025 dépasse déjà l'objectif de -40\,\% fixé pour 2026.}
\label{fig:blessures}
\end{figure}

L'objectif stratégique est de réduire le taux de fréquence de blessures de \textbf{40\,\%} d'ici 2026 par rapport à 2019 (base 2,951 blessures par million d'heures). En 2025, la réduction atteint déjà \textbf{51\,\%}, ce qui dépasse l'ambition fixée \cite{iveco_annual_2025}. Cette performance traduit un renforcement culturel, notamment via des campagnes de sensibilisation, des outils de signalement des quasi-accidents et la participation des salariés à l'identification des risques.

\subsection{Engagement communautaire}

Iveco Group inscrit sa responsabilité sociétale autour de trois priorités~: préserver la biodiversité, réduire les inégalités et protéger les groupes vulnérables, favoriser la santé et le bien-être \cite{iveco_sustainability_inaction_2025}. Les actions sont mesurées selon le cadre \glos{B4SI} (Business for Societal Impact), avec une assurance tierce pour la troisième année consécutive \cite{iveco_community_2025}.

En 2025, le groupe a consacré \textbf{1\,360\,997~\texteuro} aux initiatives communautaires, dont 91\,\% de dons en numéraire, 5\,\% en nature, 2\,\% de contributions en temps et 2\,\% de frais de gestion. Les bénéficiaires directs s'élèvent à \textbf{149\,690} personnes, en partenariat avec \textbf{139} organisations, soit un impact positif pour plus de \textbf{150\,000} personnes \cite{iveco_community_2025}.

La répartition thématique met en avant l'éducation (53\,\%), la santé (20\,\%) et l'action sociale (8\,\%). Le projet \textit{Coding the Future: Girls in Tech} illustre l'ouverture du groupe à l'inclusion numérique, tandis que le bénévolat d'entreprise mobilise les compétences métier (y compris informatiques) au profit de structures associatives \cite{iveco_community_2025, iveco_local_corp_2025}.

\begin{figure}[htbp]
\centering
\begin{tikzpicture}
\begin{axis}[
    ybar,
    width=0.85\textwidth,
    height=0.32\textwidth,
    ylabel={Part du budget communautaire (\%)},
    ymin=0, ymax=100,
    xtick=data,
    symbolic x coords={Numéraire, Nature, Temps, Gestion},
    bar width=22pt,
    nodes near coords,
    enlarge x limits=0.25,
]
\addplot[fill=ivecoblue] coordinates {(Numéraire,91) (Nature,5) (Temps,2) (Gestion,2)};
\end{axis}
\end{tikzpicture}
\caption{Répartition du budget communautaire d'Iveco Group en 2025. Les dons en numéraire dominent (91\,\%), complétés par des contributions en nature, en temps et par les frais de gestion.}
\label{fig:budgetcom}
\end{figure}

\subsection{Performance et limites}

Cette partie sociale démontre une mise en \oe{}uvre structurée~: les certifications (EDGE, ISO 45001), les données chiffrées et la mobilisation des salariés en attestent. Les axes de vigilance demeurent la progression de la part féminine dans l'effectif global, la continuité de l'équité de rémunération, et le passage d'une mesure des activités communautaires à une évaluation fine de leur impact transformationnel sur les bénéficiaires, comme le souligne l'assurance B4SI \cite{iveco_community_2025}.

\subsection*{Avis et regard critique}

La politique sociale d'Iveco Group est globalement convaincante sur le papier~: certifications reconnues, objectifs chiffrés et indicateurs audités. Sur le plan du numérique responsable, les actions communautaires (\textit{Coding the Future}) et le bénévolat de compétences offrent des points d'ancrage concrets pour l'équipe de Haute-Goulaine. Cependant, les données sur le handicap restent partielles en l'absence d'obligations légales dans certains pays, ce qui limite la comparabilité interne. Enfin, la faible part de frais de gestion (2\,\%) du budget communautaire est positive, mais l'efficacité transformative des actions sur les bénéficiaires mériterait une évaluation plus approfondie.
